% GENERATED FILE. Edit canonical lessons, then run npm run generate:books.
% canonical-source-hash: fnv1a64:f75a9cadffe86d4e
% canonical-lessons: MR-C293-bhukamp, MR-C293-hadarne, MR-C293-kosalne, MR-C293-khacne, MR-C293-darad

\chapter{When the Earth Shook}
\label{ch:mr-when-the-earth-shook}
\hlchaptermodality{293}

\begin{chapteropening}
\textbf{By the end of this chapter:} \emph{I can describe an earthquake: the ground shakes, walls fall, roads cave in.}

Complete the last lesson of chapter 293: I can describe an earthquake: the ground shakes, walls fall, roads cave in.
\end{chapteropening}

\section[\texorpdfstring{\mr{भूकंप} (bhūkamp)}{भूकंप (bhūkamp)}]{\mr{भूकंप} (bhūkamp) — an earthquake}

\label{lesson:MR-C293-bhukamp}



\begin{quote}
Before the new one: say the Marathi for ash, then the Marathi for a fire brigade.
\end{quote}

\subsection*{You'll want to know: \mr{भूकंप}}
\textbf{\mr{भूकंप}} — \emph{bhūkamp} — \textquotedblleft{}an earthquake\textquotedblright{}.

\begin{quote}
\mr{जर} \mr{भूकंप} \mr{झाला} \mr{तर} \mr{आम्ही} \mr{मोकळ्या} \mr{मैदानात} \mr{जाऊ}.
\end{quote}

\begin{quote}
\emph{jar bhūkamp jhālā tar āmhī mokaḷyā maidānāt jāū.}
\end{quote}

\begin{quote}
If there is an earthquake, we will go out to open ground.
\end{quote}

Masculine.

\mr{भूकंप} \mr{होणे}, an earthquake happens; \mr{जर} … \mr{झाला} \mr{तर}, if one happens, with the past after \mr{जर} and the future after \mr{तर}.

\begin{cousinweb}
\textbf{\mr{भूकंप}} (\emph{bhūkamp}) is a tatsama, taken from Sanskrit \emph{bhūkampa}, 'an earthquake' (Monier-Williams): \emph{bhū}, 'the earth', and \emph{kampa}, 'a trembling'.
\end{cousinweb}

\begin{grammarlens}[title={What you've built}]
One more word for telling what happened.
\end{grammarlens}

\subsection*{Your turn}
\begin{itemize}
\raggedright
  \item \emph{Say it:} \emph{bhūkamp}
  \item \emph{Say it:} \emph{bhūkamp}, once more
  \item \emph{Say it:} \emph{jar bhūkamp jhālā tar āmhī mokaḷyā maidānāt jāū.} — If there is an earthquake, we will go out to open ground.
  \item \emph{Recall:} say the Marathi for ash, then the Marathi for a fire brigade, then say \emph{bhūkamp} again
\end{itemize}

\begin{tcolorbox}[breakable,colback=teal!4,colframe=teal!35!black,title={Before you move on}]
What does \mr{भूकंप} mean? (\textquotedblleft{}An earthquake\textquotedblright{}.) Say it once more.
\end{tcolorbox}
\section[\texorpdfstring{\mr{हादरणे} (hādarṇe)}{हादरणे (hādarṇe)}]{\mr{हादरणे} (hādarṇe) — to shake; to be shaken}

\label{lesson:MR-C293-hadarne}



\begin{quote}
Before the new one: say the Marathi for a fire brigade, then the Marathi for an earthquake.
\end{quote}

\subsection*{You'll want to know: \mr{हादरणे}}
\textbf{\mr{हादरणे}} — \emph{hādarṇe} — \textquotedblleft{}to shake; to be shaken\textquotedblright{}.

\begin{quote}
\mr{भूकंपात} \mr{आमचं} \mr{संपूर्ण} \mr{घर} \mr{हादरलं}.
\end{quote}

\begin{quote}
\emph{bhūkampāt āmcaṁ sampūrṇa ghar hādarlaṁ.}
\end{quote}

\begin{quote}
Our whole house shook in the earthquake.
\end{quote}

Past (he …): \mr{तो} \mr{हादरला}, \emph{to hādarlā}.

Also of people: \mr{ती} \mr{बातमी} \mr{ऐकून} \mr{तो} \mr{हादरला}, the news shook him.

\begin{cousinweb}
\textbf{\mr{हादरणे}} (\emph{hādarṇe}) is a Marathi verb of violent shaking; its origin is not settled.
\end{cousinweb}

\begin{grammarlens}[title={What you've built}]
One more word for telling what happened.
\end{grammarlens}

\subsection*{Your turn}
\begin{itemize}
\raggedright
  \item \emph{Say it:} \emph{hādarṇe}
  \item \emph{Say it:} \emph{hādarṇe}, once more
  \item \emph{Say it:} \emph{bhūkampāt āmcaṁ sampūrṇa ghar hādarlaṁ.} — Our whole house shook in the earthquake.
  \item \emph{Recall:} say the Marathi for a fire brigade, then the Marathi for an earthquake, then say \emph{hādarṇe} again
\end{itemize}

\begin{tcolorbox}[breakable,colback=teal!4,colframe=teal!35!black,title={Before you move on}]
What does \mr{हादरणे} mean? (\textquotedblleft{}To shake; to be shaken\textquotedblright{}.) Say it once more.
\end{tcolorbox}
\section[\texorpdfstring{\mr{कोसळणे} (kosaḷṇe)}{कोसळणे (kosaḷṇe)}]{\mr{कोसळणे} (kosaḷṇe) — to collapse; to come crashing down}

\label{lesson:MR-C293-kosalne}



\begin{quote}
Before the new one: say the Marathi for an earthquake, then the Marathi for to shake.
\end{quote}

\subsection*{You'll want to know: \mr{कोसळणे}}
\textbf{\mr{कोसळणे}} — \emph{kosaḷṇe} — \textquotedblleft{}to collapse; to come crashing down\textquotedblright{}.

\begin{quote}
\mr{जुनी} \mr{भिंत} \mr{पावसात} \mr{कोसळली}.
\end{quote}

\begin{quote}
\emph{junī bhint pāvsāt kosaḷlī.}
\end{quote}

\begin{quote}
The old wall collapsed in the rain.
\end{quote}

Past (he …): \mr{तो} \mr{कोसळला}, \emph{to kosaḷlā}.

Also of rain: \mr{पाऊस} \mr{कोसळतो}, the rain pours down.

\begin{cousinweb}
\textbf{\mr{कोसळणे}} (\emph{kosaḷṇe}) is a Marathi verb of falling in a heap; its origin is not settled.
\end{cousinweb}

\begin{grammarlens}[title={What you've built}]
One more word for telling what happened.
\end{grammarlens}

\subsection*{Your turn}
\begin{itemize}
\raggedright
  \item \emph{Say it:} \emph{kosaḷṇe}
  \item \emph{Say it:} \emph{kosaḷṇe}, once more
  \item \emph{Say it:} \emph{junī bhint pāvsāt kosaḷlī.} — The old wall collapsed in the rain.
  \item \emph{Recall:} say the Marathi for an earthquake, then the Marathi for to shake, then say \emph{kosaḷṇe} again
\end{itemize}

\begin{tcolorbox}[breakable,colback=teal!4,colframe=teal!35!black,title={Before you move on}]
What does \mr{कोसळणे} mean? (\textquotedblleft{}To collapse; to come crashing down\textquotedblright{}.) Say it once more.
\end{tcolorbox}
\section[\texorpdfstring{\mr{खचणे} (khacṇe)}{खचणे (khacṇe)}]{\mr{खचणे} (khacṇe) — to cave in; to sink down}

\label{lesson:MR-C293-khacne}



\begin{quote}
Before the new one: say the Marathi for to shake, then the Marathi for to collapse.
\end{quote}

\subsection*{You'll want to know: \mr{खचणे}}
\textbf{\mr{खचणे}} — \emph{khacṇe} — \textquotedblleft{}to cave in; to sink down\textquotedblright{}.

\begin{quote}
\mr{पावसामुळे} \mr{गावातला} \mr{रस्ता} \mr{खचला}.
\end{quote}

\begin{quote}
\emph{pāvsāmuḷe gāvātlā rastā khaclā.}
\end{quote}

\begin{quote}
Because of the rain, the village road caved in.
\end{quote}

Past (he …): \mr{तो} \mr{खचला}, \emph{to khaclā}.

Also of courage: \mr{तो} \mr{मनाने} \mr{खचला}, he lost heart.

\begin{cousinweb}
\textbf{\mr{खचणे}} (\emph{khacṇe}) is inherited: CDIAL lists Marathi \emph{khacṇe}, 'to sink down', under \emph{*khacc}, 'sink down or back' (CDIAL 3764).
\end{cousinweb}

\begin{grammarlens}[title={What you've built}]
One more word for telling what happened.
\end{grammarlens}

\subsection*{Your turn}
\begin{itemize}
\raggedright
  \item \emph{Say it:} \emph{khacṇe}
  \item \emph{Say it:} \emph{khacṇe}, once more
  \item \emph{Say it:} \emph{pāvsāmuḷe gāvātlā rastā khaclā.} — Because of the rain, the village road caved in.
  \item \emph{Recall:} say the Marathi for to shake, then the Marathi for to collapse, then say \emph{khacṇe} again
\end{itemize}

\begin{tcolorbox}[breakable,colback=teal!4,colframe=teal!35!black,title={Before you move on}]
What does \mr{खचणे} mean? (\textquotedblleft{}To cave in; to sink down\textquotedblright{}.) Say it once more.
\end{tcolorbox}
\section[\texorpdfstring{\mr{दरड} (daraḍ)}{दरड (daraḍ)}]{\mr{दरड} (daraḍ) — a landslide}

\label{lesson:MR-C293-darad}



\begin{quote}
Before the new one: say the Marathi for to collapse, then the Marathi for to cave in.
\end{quote}

\subsection*{You'll want to know: \mr{दरड}}
\textbf{\mr{दरड}} — \emph{daraḍ} — \textquotedblleft{}a landslide\textquotedblright{}.

\begin{quote}
\mr{घाटात} \mr{दरड} \mr{कोसळली} \mr{आणि} \mr{वाहतूक} \mr{थांबली}.
\end{quote}

\begin{quote}
\emph{ghāṭāt daraḍ kosaḷlī āṇi vāhtūk thāmblī.}
\end{quote}

\begin{quote}
A landslide came down in the pass and the traffic stopped.
\end{quote}

Feminine.

\mr{दरड} \mr{कोसळणे}, a landslide comes down.

\begin{cousinweb}
\textbf{\mr{दरड}} (\emph{daraḍ}) is inherited: CDIAL lists Marathi \emph{daraḍ}, 'a steep bank', under Sanskrit \emph{darad}, 'a bank, a precipice' (CDIAL 6191).
\end{cousinweb}

\begin{grammarlens}[title={What you've built}]
One more word for telling what happened.
\end{grammarlens}

\subsection*{Your turn}
\begin{itemize}
\raggedright
  \item \emph{Say it:} \emph{daraḍ}
  \item \emph{Say it:} \emph{daraḍ}, once more
  \item \emph{Say it:} \emph{ghāṭāt daraḍ kosaḷlī āṇi vāhtūk thāmblī.} — A landslide came down in the pass and the traffic stopped.
  \item \emph{Recall:} say the Marathi for to collapse, then the Marathi for to cave in, then say \emph{daraḍ} again
\end{itemize}

\begin{tcolorbox}[breakable,colback=teal!4,colframe=teal!35!black,title={Before you move on}]
What does \mr{दरड} mean? (\textquotedblleft{}A landslide\textquotedblright{}.) Say it once more.
\end{tcolorbox}
